\documentclass[11pt,a4paper]{article}

\usepackage[T1]{fontenc}
\usepackage[utf8]{inputenc}
\usepackage{lmodern}
\usepackage{microtype}
\usepackage[a4paper,margin=27mm,headheight=14pt]{geometry}
\usepackage{amsmath,amssymb,amsthm,mathtools}
\usepackage{booktabs,longtable,array,tabularx}
\usepackage{enumitem}
\usepackage{xcolor}
\usepackage{hyperref}
\usepackage{fancyhdr}
\usepackage{listings}
\usepackage{upquote}
\usepackage{parskip}

\definecolor{leanblue}{RGB}{20,62,116}
\definecolor{darkgray}{RGB}{70,70,70}
\definecolor{lightgray}{RGB}{246,247,249}
\definecolor{darkred}{RGB}{128,20,20}
\definecolor{darkgreen}{RGB}{15,95,55}

\hypersetup{
  colorlinks=true,
  linkcolor=leanblue,
  citecolor=darkgreen,
  urlcolor=leanblue,
  pdftitle={Planar negative-power circular maximal theorem: formalization blueprint},
  pdfauthor={Prepared for LeanSpherical}
}

\pagestyle{fancy}
\fancyhf{}
\lhead{Duoandikoetxea--Vega planar blueprint}
\rhead{LeanSpherical}
\cfoot{\thepage}

\lstdefinelanguage{Lean}{
  morekeywords={theorem,lemma,def,abbrev,namespace,end,open,scoped,notation,section,variable,include,by,exact,apply,refine,have,show,from,let,in,if,then,else,match,with,where,fun,forall,exists,import,noncomputable,instance,structure,class,extends},
  sensitive=true,
  morecomment=[l]{--},
  morecomment=[s]{/-}{-/},
  morestring=[b]"
}

\lstset{
  language=Lean,
  basicstyle=\ttfamily\small,
  keywordstyle=\color{leanblue}\bfseries,
  commentstyle=\color{darkgreen},
  stringstyle=\color{darkred},
  backgroundcolor=\color{lightgray},
  frame=single,
  rulecolor=\color{gray!35},
  breaklines=true,
  breakatwhitespace=false,
  columns=fullflexible,
  keepspaces=true,
  showstringspaces=false,
  tabsize=2,
  xleftmargin=3mm,
  xrightmargin=3mm
}

\newtheorem{theorem}{Theorem}[section]
\newtheorem{proposition}[theorem]{Proposition}
\newtheorem{lemma}[theorem]{Lemma}
\newtheorem{corollary}[theorem]{Corollary}
\newtheorem{claim}[theorem]{Claim}
\theoremstyle{definition}
\newtheorem{definition}[theorem]{Definition}
\newtheorem{blueprint}[theorem]{Blueprint item}
\theoremstyle{remark}
\newtheorem{remark}[theorem]{Remark}
\newtheorem{warning}[theorem]{Formalization warning}

\newcommand{\R}{\mathbb{R}}
\newcommand{\C}{\mathbb{C}}
\newcommand{\N}{\mathbb{N}}
\newcommand{\E}{\mathcal{E}}
\newcommand{\M}{\mathcal{M}}
\newcommand{\Mc}{\mathcal{M}_{\mathrm{c}}}
\newcommand{\A}{\mathcal{A}}
\newcommand{\Lp}{L^{p}}
\newcommand{\walpha}{w_{\alpha}}
\newcommand{\norm}[1]{\left\lVert #1\right\rVert}
\newcommand{\abs}[1]{\left\lvert #1\right\rvert}
\newcommand{\one}{\mathbf{1}}
\newcommand{\esssup}{\mathop{\mathrm{ess\,sup}}}
\newcommand{\into}{\longrightarrow}
\newcommand{\lean}{\texttt}

\title{\textbf{The planar negative-power circular maximal theorem}\\[4pt]
\large A Duoandikoetxea--Vega blueprint and the exact completion of\\
\lean{PowerWeights.closure\_typeSet\_eq} in dimension two}
\author{Prepared for the \texttt{LeanSpherical} formalization}
\date{25 August 2026}

\begin{document}
\maketitle

\begin{abstract}
The public theorem in \texttt{LeanSpherical.Theorems} currently assumes
$d\geq 3$.  The only additional analytic input needed for $d=2$ is a
negative-power weighted circular maximal estimate.  The direct drop-in form is
boundedness of the full circular maximal operator on
$L^p(\R^2,\abs{x}^{\alpha}\,dx)$ for $p>2$ and $-1<\alpha<0$.  This is a
strict subrange of Theorem~10 of Duoandikoetxea and Vega, whose planar range is
$-1<\alpha<p-2$.

This document gives: the exact theorem statements; a complete reduction of the
$d=2$ Fraccaroli--Roos--Seeger result to that input; a source audit; a
formalization-oriented decomposition of the Duoandikoetxea--Vega proof; proposed
Lean declarations; and a dependency and risk checklist.  It distinguishes the
minimal local theorem needed by the existing FRS globalization lemma from the
more convenient global theorem that can be exposed as the public API.
\end{abstract}

\tableofcontents
\newpage

\section{Executive conclusion}

\subsection{The exact missing result}

Let $\sigma$ be normalized arclength measure on the unit circle $S^1$, and put

\[
\A_t f(x)=\int_{S^1} f(x+t\omega)\,d\sigma(\omega).
\]

For a set of positive radii $E\subset(0,\infty)$, put

\[
\M_E f(x)=\sup_{t\in E}\abs{\A_t f(x)}.
\]

The full circular maximal operator is

\[
\M f=\M_{(0,\infty)}f.
\]

The direct theorem needed by the public $d=2$ argument is the following.

\begin{theorem}[Planar negative-power circular maximal theorem]\label{thm:direct-missing}
Let $p$ and $\alpha$ be real numbers satisfying

\[
p>2
\]

and

\[
-1<\alpha<0.
\]

There is a constant $C_{p,\alpha}>0$ such that every measurable
$f:\R^2\to\C$ in $L^p(\R^2,\abs{x}^{\alpha}\,dx)$ satisfies

\[
\norm{\M f}_{L^p(\abs{x}^{\alpha})}
\leq C_{p,\alpha}\norm{f}_{L^p(\abs{x}^{\alpha})}.
\]
\end{theorem}

This is the cleanest theorem to import into the final branch of the FRS proof.
It immediately implies the same estimate for every restricted set of radii,
because

\[
\M_E f(x)\leq \M f(x).
\]

There is, however, a strictly smaller analytic target if the already formalized
FRS globalization lemma is used.

\begin{theorem}[Minimal local analytic input]\label{thm:minimal-local}
Let $p>2$ and $-1<\alpha<0$.  There is a constant $C_{p,\alpha}>0$ such that

\[
\norm{\M_{[1,2]} f}_{L^p(\abs{x}^{\alpha})}
\leq C_{p,\alpha}\norm{f}_{L^p(\abs{x}^{\alpha})}.
\]
\end{theorem}

Theorem~\ref{thm:minimal-local}, Bourgain's unweighted circular maximal theorem,
and FRS Lemma~2.1 imply Theorem~\ref{thm:direct-missing} and, more generally,
the required bound for every $E\subset(0,\infty)$.  Thus:

\begin{itemize}[leftmargin=7mm]
\item Theorem~\ref{thm:minimal-local} is the minimal new analytic theorem.
\item Theorem~\ref{thm:direct-missing} is the best direct Lean interface.
\item The full Duoandikoetxea--Vega theorem is stronger and is the natural
      source theorem to formalize.
\end{itemize}

\subsection{The source theorem}

Duoandikoetxea and Vega prove the following result in arbitrary dimension.

\begin{theorem}[Duoandikoetxea--Vega, Theorem 10]\label{thm:dv-general}
Let $d\geq2$.  If

\[
p>\frac{d}{d-1}
\]

and

\[
1-d<\alpha<(d-1)(p-1)-1,
\]

then the full spherical maximal operator is bounded on
$L^p(\R^d,\abs{x}^{\alpha}\,dx)$.
\end{theorem}

For $d=2$, the range becomes

\[
p>2
\]

and

\[
-1<\alpha<p-2.
\]

Consequently Theorem~\ref{thm:direct-missing} is an immediate planar subcase:
if $p>2$ and $\alpha<0$, then

\[
\alpha<0<p-2.
\]

\subsection{Verdict on the rest of the $d=2$ proof}

After Theorem~\ref{thm:direct-missing} or
Theorem~\ref{thm:minimal-local} is available, no additional planar harmonic
analysis is needed for the closure theorem.  The existing proof divides into
three exhaustive finite-$p$ branches:

\begin{enumerate}[leftmargin=8mm]
\item $\alpha\geq0$: the already formalized nonnegative-weight argument;
\item $\alpha<0$ and $p\leq2$: the already formalized FRS argument;
\item $\alpha<0$ and $p>2$: the new Duoandikoetxea--Vega branch.
\end{enumerate}

The necessity argument, convexity of the type set, the passage from strict
inequalities to the closed admissible region, and the vertical $1/p=0$ segment
are unchanged.  The endpoint $\alpha=-1$ is not needed as an analytic theorem:
when it belongs to the admissible boundary it is recovered by closure.

\section{Conventions and comparison with the Lean definitions}

\subsection{Averages and signs}

The project uses

\[
\operatorname{sphericalAverage}(t,f,x)
=
\int_{S^{d-1}} f(x+t\omega)\,d\sigma(\omega).
\]

Some sources write $f(x-t\omega)$.  The two conventions agree because the map
$\omega\mapsto-\omega$ preserves spherical measure.  This should be made a
named lemma rather than hidden in a change of variables.

The project normalizes surface measure.  Multiplying surface measure by a fixed
positive scalar only changes the operator norm, so source results stated with
unnormalized surface measure transfer immediately after one scalar calculation.

\subsection{Power weights}

The project defines

\[
\operatorname{powerWeight}(d,\alpha)
=
\abs{x}^{\alpha}\,dx
\]

through \lean{volume.withDensity}.  For $\alpha<0$, the density is infinite at
$x=0$ in the \lean{ENNReal} model.  This is harmless because $\{0\}$ is
Lebesgue-null, but identities involving the density must be stated almost
everywhere.  Pointwise rewrites at the origin are the wrong interface.

\subsection{Maximal-function codomain}

The project maximal function is \lean{ENNReal}-valued:

\[
M(E,f,x)=\sup_{t\in E\cap(0,\infty)}
\operatorname{ofReal}\bigl(\norm{\A_t f(x)}\bigr).
\]

The analytic literature usually uses a real-valued nonnegative function.  A
small adapter layer should establish the almost-everywhere equality between the
project expression and \lean{ENNReal.ofReal} of the classical supremum.  Once
that layer exists, all norm inequalities should be proved at one codomain and
transported once, rather than repeatedly unfolding \lean{M}.

\section{The complete $d=2$ reduction}

This section proves that the new planar theorem is exactly sufficient for the
public closure theorem.  No analytic estimate from Duoandikoetxea--Vega is used
until the final line of Proposition~\ref{prop:high-p-branch}.

\subsection{The strict FRS inequalities}

For the sufficiency part of the FRS theorem, one proves boundedness at points
satisfying strict versions of the defining inequalities.  In dimension two the
lower strict inequality can be written in the form

\[
\nu_E^\sharp(p-2-\alpha)<p-1.
\]

The FRS theory also gives the elementary lower bound

\[
\nu_E^\sharp(\rho)\geq\rho
\]

for every real $\rho$.  These two facts force the lower Duoandikoetxea--Vega
weight condition.

\begin{lemma}[The planar strict region lies above $\alpha=-1$]
Assume

\[
\nu_E^\sharp(p-2-\alpha)<p-1.
\]

Then

\[
\alpha>-1.
\]
\end{lemma}

\begin{proof}
The general inequality for the Legendre--Assouad function gives

\[
p-2-\alpha\leq\nu_E^\sharp(p-2-\alpha).
\]

Combining it with the hypothesis gives

\[
p-2-\alpha<p-1.
\]

Subtracting $p-2$ from both sides gives

\[
-\alpha<1.
\]

Multiplying by $-1$ reverses the inequality and gives

\[
\alpha>-1.
\]
\end{proof}

\begin{lemma}[The high-$p$ negative branch lies below $p-2$]
If $p>2$ and $\alpha<0$, then

\[
\alpha<p-2.
\]
\end{lemma}

\begin{proof}
The first hypothesis gives

\[
0<p-2.
\]

The second gives

\[
\alpha<0.
\]

Transitivity yields

\[
\alpha<p-2.
\]
\end{proof}

\subsection{The only new branch}

\begin{proposition}[Completion of the finite-$p$ sufficiency proof]\label{prop:high-p-branch}
Assume the strict FRS hypotheses in dimension two.  Assume also that
$p>2$ and $\alpha<0$.  Then $\M_E$ is bounded on
$L^p(\R^2,\abs{x}^{\alpha}\,dx)$.
\end{proposition}

\begin{proof}
The strict lower FRS inequality and the preceding lemma imply

\[
-1<\alpha.
\]

The hypotheses $p>2$ and $\alpha<0$ imply

\[
\alpha<p-2.
\]

Therefore the planar case of Duoandikoetxea--Vega Theorem~10 applies and gives

\[
\norm{\M f}_{L^p(\abs{x}^{\alpha})}
\leq C_{p,\alpha}\norm{f}_{L^p(\abs{x}^{\alpha})}.
\]

Since $E\subset(0,\infty)$, one has pointwise

\[
\M_E f\leq\M f.
\]

Monotonicity of the $L^p$ norm gives

\[
\norm{\M_E f}_{L^p(\abs{x}^{\alpha})}
\leq C_{p,\alpha}\norm{f}_{L^p(\abs{x}^{\alpha})}.
\]
\end{proof}

\subsection{A proof using only the minimal local theorem}

Suppose only Theorem~\ref{thm:minimal-local} has been formalized.  For $R>0$,
define the normalized slice

\[
E_R=R^{-1}\bigl(E\cap[R,2R]\bigr).
\]

Then

\[
E_R\subset[1,2].
\]

Consequently

\[
\M_{E_R}f\leq\M_{[1,2]}f.
\]

The local theorem bounds the second term in the FRS globalization estimate,
uniformly in $R$.  Bourgain's circular maximal theorem and
$\M_E\leq\M$ bound the unweighted term.  FRS Lemma~2.1 then gives the desired
global weighted estimate.  This route proves the restricted theorem directly
and avoids first constructing a standalone full-radii weighted theorem.

\subsection{Passage to the closure}

The analytic estimates prove membership in \lean{typeSet} only under strict
inequalities.  The existing FRS argument then proceeds as follows.

\begin{enumerate}[leftmargin=8mm]
\item Approximate every finite-$p$ boundary point of the admissible region by
      strict interior points.
\item Use the definition of topological closure to place the boundary point in
      \lean{closure (typeSet 2 E)}.
\item Approximate the vertical segment with first coordinate zero by points
      with $p\to\infty$.
\item Combine with the already proved necessary inclusion

\[
\operatorname{closure}(\operatorname{typeSet}(2,E))
\subseteq
\operatorname{admissibleRegion}(2,E).
\]
\end{enumerate}

The new theorem affects only item 1 for strict points with $p>2$ and
$\alpha<0$.  It does not alter any closure or convexity argument.

\section{Source audit: what Duoandikoetxea--Vega actually prove}

\subsection{Bibliographic identification}

The relevant source is:

\begin{quote}
Javier Duoandikoetxea and Luis Vega, ``Spherical means and weighted
inequalities,'' \emph{Journal of the London Mathematical Society} (2),
53(2):343--353, 1996.  DOI: 10.1112/jlms/53.2.343.
\end{quote}

The required result is Theorem~10.  The source statement is precisely
Theorem~\ref{thm:dv-general}.  In the proof, the authors dispatch the case
$\alpha\geq0$ by their Proposition~8.  For $\alpha<0$, they reduce to a
unit-scale estimate, displayed as their equation (2), and use a
frequency-localized estimate displayed as equation (3).  A later endpoint
paper by Juyoung Lee explicitly identifies equation (3) in the proof of
Theorem~10 and records that it carries a $2^{j\varepsilon}$ loss; Lee also says
that one of the two spatial-frequency regimes is essentially the argument of
Duoandikoetxea--Vega.

\subsection{Access and reconstruction note}

The publisher page exposes the exact theorem statement and the beginning of the
proof through indexed full text, but the publisher PDF was not downloadable in
the research environment used for this blueprint.  The source-level structure
above was therefore checked in three ways:

\begin{enumerate}[leftmargin=8mm]
\item the indexed text of Theorem~10 and its first proof reduction;
\item the bibliographic and range statement in the FRS paper;
\item Lee's page-level discussion of equation (3), the $2^{j\varepsilon}$
      loss, and the annular argument inherited from the 1996 paper.
\end{enumerate}

The blueprint below is a formalization-oriented reconstruction of that proof
architecture.  It is not presented as a line-by-line transcription of the
publisher typesetting.  Every theorem used in the reconstruction is stated
explicitly, so the formal development does not depend on unrecorded source
notation.

\subsection{What should and should not be formalized}

The 1996 paper contains more than the missing planar branch.  For the present
project it is unnecessary to formalize:

\begin{itemize}[leftmargin=7mm]
\item the entire dyadic-radii theory developed earlier in the paper;
\item the nonnegative-weight branch, already available in the project;
\item the arbitrary-dimensional theorem before completing the planar branch;
\item the endpoint $\alpha=-1$.
\end{itemize}

The right initial target is the local planar theorem for $p>2$ and
$-1<\alpha<0$, followed by the global and restricted corollaries.

\section{Analytic proof blueprint}

\subsection{Overview}

The proof is organized around a Littlewood--Paley decomposition.  The
unweighted local circular maximal estimate has a decaying frequency norm for
$p>2$.  At the critical weight $\abs{x}^{-1}$ one has a frequency estimate
with an arbitrarily small exponential loss.  Interpolation with change of
measure produces exponential decay at every intermediate weight
$\abs{x}^{\alpha}$ with $-1<\alpha<0$.  The decay makes the frequency sum
converge.

The following items are stated for $d=2$, because that is the only new case.
They are deliberately split at interfaces that can be tested independently in
Lean.

\subsection{Phase A: the local operator and a countable supremum}

\begin{blueprint}[Local operator]
Define

\[
\Mc f(x)=\sup_{1\leq t\leq2}\abs{\A_t f(x)}.
\]

For Schwartz $f$, prove that $t\mapsto\A_t f(x)$ is continuous for every $x$.
If $D=\mathbb{Q}\cap[1,2]$, deduce

\[
\Mc f(x)=\sup_{t\in D}\abs{\A_t f(x)}.
\]

This converts the continuum supremum to a countable supremum and gives
measurability without a separate measurable-maximum theorem.
\end{blueprint}

\begin{blueprint}[Finite-grid approximants]
Choose increasing finite sets $D_N\subset D$ whose union is dense in $[1,2]$,
and define

\[
\mathcal{M}_{\mathrm{c},N} f(x)=\max_{t\in D_N}\abs{\A_t f(x)}.
\]

Then

\[
\mathcal{M}_{\mathrm{c},N} f(x)\uparrow\Mc f(x).
\]

All interpolation arguments should first be carried out uniformly in $N$ and
then passed to the limit by monotone convergence or Fatou's lemma.  This avoids
formalizing complex interpolation for a nonlinear uncountable supremum.
\end{blueprint}

\subsection{Phase B: Littlewood--Paley pieces}

Choose a smooth radial cutoff $\chi$ supported in an annulus and satisfying a
partition of unity away from the origin.  Let $P_j$ be the Fourier multiplier
with symbol $\chi(2^{-j}\abs{\xi})$.  Define

\[
T_j f(x,t)=\A_t(P_jf)(x)
\]

and

\[
M_j f(x)=\sup_{1\leq t\leq2}\abs{T_jf(x,t)}.
\]

Use a separate smooth low-frequency projection $P_{\mathrm{lo}}$.

\begin{blueprint}[Reconstruction inequality]
For Schwartz $f$, prove the pointwise estimate

\[
\Mc f(x)
\leq
\Mc(P_{\mathrm{lo}}f)(x)
+
\sum_{j\geq1}M_jf(x).
\]

It is enough to use the triangle inequality after reconstructing $f$ from the
frequency partition.  The formal statement should be first proved for finite
partial sums and then passed to the limit.
\end{blueprint}

\begin{blueprint}[Low frequencies]
Prove

\[
\Mc(P_{\mathrm{lo}}f)(x)
\leq C\,M_{\mathrm{HL}}f(x),
\]

where $M_{\mathrm{HL}}$ is the Hardy--Littlewood maximal function.  Hence the
low-frequency term is bounded on $L^p(\abs{x}^{\alpha})$ whenever
$\abs{x}^{\alpha}$ is an $A_p$ weight.
\end{blueprint}

For $d=2$, the power-weight criterion is

\[
\abs{x}^{\alpha}\in A_p
\quad\Longleftrightarrow\quad
-2<\alpha<2(p-1).
\]

The desired range $p>2$ and $-1<\alpha<0$ lies strictly inside this interval.
If a general $A_p$ library is not available, only the special power-weight
case and the convolution kernels arising here need to be proved.

\subsection{Phase C: unweighted frequency gain}

\begin{blueprint}[Bourgain frequency-localized estimate]\label{bp:bourgain-piece}
For every real $p>2$, prove that there are constants $C_p>0$ and
$\delta_p>0$ such that, for $j\geq1$,

\[
\norm{M_jf}_{L^p(\R^2)}
\leq C_p2^{-\delta_pj}\norm{f}_{L^p(\R^2)}.
\]
\end{blueprint}

This is the frequency-localized form of the local circular maximal estimate.
It is stronger data than the final statement ``$\Mc$ is bounded on $L^p$.''
A completed proof of Bourgain's theorem normally contains this estimate, but it
cannot be recovered abstractly from the final boundedness theorem.  The Lean
port should therefore check that this frequency-gain lemma is exposed by the
existing Bourgain formalization.

For finite grids $D_N$, use the linear map

\[
f\longmapsto \bigl(T_jf(\,\boldsymbol{\cdot}\,,t)\bigr)_{t\in D_N}
\]

with target $L^p(\R^2;\ell^\infty(D_N))$.  The estimate must be uniform in $N$.

\subsection{Phase D: the critical-weight estimate with loss}

The hard source estimate is at the critical weight $\abs{x}^{-1}$.  The form
best suited to interpolation is the following packaged proposition.

\begin{proposition}[Critical-weight dyadic estimate]\label{prop:critical-loss}
For every $\varepsilon>0$, there is $C_\varepsilon>0$ such that, for $j\geq1$,

\[
\norm{M_jf}_{L^2(\abs{x}^{-1})}
\leq C_\varepsilon 2^{\varepsilon j}
\norm{f}_{L^2(\abs{x}^{-1})}.
\]

The same estimate holds for finite-grid maximal operators with a constant
independent of the grid.
\end{proposition}

Proposition~\ref{prop:critical-loss} is the formalization package obtained from
the annular estimate in equation (3) of the proof of Duoandikoetxea--Vega
Theorem~10 and the subsequent annular summation.  The $2^{\varepsilon j}$ loss
is essential to the source interpolation argument; it is harmless away from
$\alpha=-1$.

A source-faithful proof should be divided as follows.

\begin{blueprint}[Spatial annuli]
Let

\[
\mathcal{R}_k=
\left\{x\in\R^2:2^{-k-1}<\abs{x}\leq2^{-k}\right\}.
\]

Decompose both the input and output into dyadic annuli.  On $\mathcal{R}_k$ the
weight is comparable to the constant $2^k$:

\[
\abs{x}^{-1}\asymp2^k.
\]

Record the comparison as two explicit inequalities, because Lean will not use
asymptotic notation.
\end{blueprint}

\begin{blueprint}[Near and far interactions]
For fixed frequency $2^j$, split the annular interactions according to the
relative sizes of $j$ and $k$.  The geometric regime in which the spatial
annulus is smaller than the oscillation scale is estimated by the source
annular argument.  Off-diagonal interactions are controlled by rapid decay of
the Schwartz kernels.  State each regime with a numerical decay factor and a
summable majorant before performing any double sum.
\end{blueprint}

\begin{blueprint}[Maximal parameter]
Use a one-dimensional Sobolev estimate in $t$ on $[1,2]$:

\[
\sup_{1\leq t\leq2}\abs{F(t)}
\leq
C\norm{F}_{L^2([1,2])}^{1/2}
\norm{F'}_{L^2([1,2])}^{1/2}
+
C\norm{F}_{L^2([1,2])}.
\]

Apply this with $F(t)=\A_t(P_jf)(x)$.  Frequency localization controls the
$t$-derivative.  Keep the averaging estimate and the derivative estimate as
separate lemmas; this is the point at which powers of $2^j$ are generated.
\end{blueprint}

\begin{blueprint}[Annular summation]
After proving the localized bounds, sum first in the off-diagonal annulus
index, then in the main index.  Introduce a small parameter $\varepsilon>0$ to
majorize logarithmic or borderline summation by $2^{\varepsilon j}$.  The final
result is Proposition~\ref{prop:critical-loss}.
\end{blueprint}

\begin{remark}
Juyoung Lee's later endpoint paper proves sharper annular estimates and removes
the loss in the global endpoint theorem for the admissible values of $p$.  That
stronger endpoint theory is not needed here.  The lossy source estimate is
exactly what interpolation for $-1<\alpha<0$ requires.
\end{remark}

\subsection{Phase E: move the critical estimate from $p=2$ to $p>2$}

For every $j$ one has the trivial bound

\[
\norm{M_jf}_{L^\infty}
\leq C\norm{f}_{L^\infty}.
\]

The same statement is valid when both spaces are regarded over the fixed
measure $\abs{x}^{-1}\,dx$, because $L^\infty$ does not depend on multiplication
of the measure by a positive density up to null sets.

Interpolate Proposition~\ref{prop:critical-loss} with this $L^\infty$ bound on
the fixed measure space $(\R^2,\abs{x}^{-1}\,dx)$.  If $p>2$, the interpolation
parameter satisfies

\[
\frac{1}{p}=\frac{1-\vartheta}{2}.
\]

Thus

\[
1-\vartheta=\frac{2}{p}.
\]

For every $\varepsilon>0$, after renaming the loss parameter, obtain

\[
\norm{M_jf}_{L^p(\abs{x}^{-1})}
\leq C_{p,\varepsilon}2^{\varepsilon j}
\norm{f}_{L^p(\abs{x}^{-1})}.
\]

This interpolation must again be carried out for finite-grid vector-valued
linear operators and then passed to the supremum.

\subsection{Phase F: Stein--Weiss interpolation in the weight}

Fix $p>2$ and $-1<\alpha<0$.  Put

\[
\theta=-\alpha.
\]

Then

\[
0<\theta<1.
\]

The two endpoint measures are

\[
d\mu_0=dx
\]

and

\[
d\mu_1=\abs{x}^{-1}\,dx.
\]

Stein--Weiss interpolation with change of measure gives

\[
\bigl[L^p(d\mu_0),L^p(d\mu_1)\bigr]_\theta
=
L^p(\abs{x}^{-\theta}\,dx).
\]

Since $-\theta=\alpha$, the interpolated measure is exactly
$\abs{x}^{\alpha}\,dx$.

Interpolate the unweighted estimate

\[
\norm{M_jf}_{L^p}
\leq C_p2^{-\delta_pj}\norm{f}_{L^p}
\]

with the critical-weight estimate

\[
\norm{M_jf}_{L^p(\abs{x}^{-1})}
\leq C_{p,\varepsilon}2^{\varepsilon j}
\norm{f}_{L^p(\abs{x}^{-1})}.
\]

The resulting frequency exponent is

\[
-(1-\theta)\delta_p+\theta\varepsilon.
\]

Choose $\varepsilon$ so that

\[
0<\varepsilon<\frac{1-\theta}{\theta}\delta_p.
\]

Define

\[
\eta=(1-\theta)\delta_p-\theta\varepsilon.
\]

Then

\[
\eta>0.
\]

The interpolated estimate is

\[
\norm{M_jf}_{L^p(\abs{x}^{\alpha})}
\leq C_{p,\alpha}2^{-\eta j}
\norm{f}_{L^p(\abs{x}^{\alpha})}.
\]

\subsection{Phase G: sum the frequencies}

The low-frequency term is bounded.  For the high-frequency terms, Minkowski's
inequality and the geometric series give

\[
\norm{\Mc f}_{L^p(\abs{x}^{\alpha})}
\leq
C\norm{f}_{L^p(\abs{x}^{\alpha})}
+
\sum_{j\geq1}
\norm{M_jf}_{L^p(\abs{x}^{\alpha})}.
\]

Therefore

\[
\norm{\Mc f}_{L^p(\abs{x}^{\alpha})}
\leq
C\left(1+\sum_{j\geq1}2^{-\eta j}\right)
\norm{f}_{L^p(\abs{x}^{\alpha})}.
\]

The series converges because $\eta>0$, so

\[
\norm{\Mc f}_{L^p(\abs{x}^{\alpha})}
\leq C_{p,\alpha}\norm{f}_{L^p(\abs{x}^{\alpha})}.
\]

This proves Theorem~\ref{thm:minimal-local} first for Schwartz functions.

\subsection{Phase H: extension to all $L^p$ functions}

Choose Schwartz functions $f_n$ converging to $f$ in
$L^p(\abs{x}^{\alpha})$.  The formal extension can be implemented by either of
two standard routes.

\begin{enumerate}[leftmargin=8mm]
\item Construct the bounded operator on the completion and identify it almost
      everywhere with the maximal function by a subsequence and Fatou's lemma.
\item Use sublinearity and a maximal-principle extension lemma already present
      in the project.
\end{enumerate}

The conclusion must include both \lean{MemLp} of the output and the norm
inequality, because \lean{typeSet} requires the conjunction.

\subsection{Phase I: globalization}

For the minimal implementation, apply the already formalized FRS local-to-global
lemma.  The inputs are:

\begin{itemize}[leftmargin=7mm]
\item the unweighted $L^p$ bound for $\M_E$, inherited from Bourgain by
      $\M_E\leq\M$;
\item the uniform weighted local bound for every $E_R\subset[1,2]$, inherited
      from Theorem~\ref{thm:minimal-local} by monotonicity.
\end{itemize}

The output is the weighted global bound for $\M_E$.  Taking
$E=(0,\infty)$ gives Theorem~\ref{thm:direct-missing}.

\section{Interpolation interface for Lean}

\subsection{Why the maximal operator must be linearized}

Complex interpolation theorems apply to linear or suitably analytic operator
families.  The map $f\mapsto M_jf$ is sublinear.  The FRS paper resolves the
same issue using measurable selectors and linearization.  For formalization,
finite rational grids are usually simpler.

Let $F\subset\mathbb{Q}\cap[1,2]$ be finite.  Define

\[
\mathcal{T}_{j,F}f(x)=\bigl(\A_tP_jf(x)\bigr)_{t\in F}.
\]

Give $\C^F$ the supremum norm.  Then

\[
\norm{\mathcal{T}_{j,F}f(x)}_{\ell^\infty(F)}
=
\max_{t\in F}\abs{\A_tP_jf(x)}.
\]

The map $\mathcal{T}_{j,F}$ is linear.  Prove all endpoint estimates uniformly
in $F$, interpolate the linear maps, and pass through an increasing sequence
of finite sets.  This avoids measurable selector construction and is compatible
with existing vector-valued interpolation machinery if finite products are
already supported.

\subsection{Required interpolation theorems}

The formal development needs two forms of interpolation.

\begin{enumerate}[leftmargin=8mm]
\item Fixed-measure interpolation between $L^2(\abs{x}^{-1})$ and
      $L^\infty(\abs{x}^{-1})$.
\item Stein--Weiss interpolation with change of measure at the same exponent
      $p$ between $dx$ and $\abs{x}^{-1}dx$.
\end{enumerate}

If only scalar Riesz--Thorin is available, use a finite-dimensional target and
its normed-space instance.  The change-of-measure theorem should be exposed as
a separate reusable result rather than reproved inside the circular maximal
file.

\section{Proposed Lean API}

\subsection{The minimal local theorem}

A declaration close to the current project conventions is:

\begin{lstlisting}
namespace Spherical.PowerWeights

/-- Local planar Duoandikoetxea--Vega estimate in the negative-weight range. -/
theorem eLpNorm_localCircularMaximal_powerWeight_le_of_neg
    {p : ENNReal} {alpha : Real}
    (hp_top : p != top) (hp : (2 : ENNReal) < p)
    (halpha_lo : -1 < alpha) (halpha_hi : alpha < 0) :
    exists C : Real, 0 < C and forall f : (Real^2) -> Complex,
      MemLp f p (powerWeight 2 alpha) ->
      MemLp (M (Icc (1 : Real) 2) f) p (powerWeight 2 alpha) and
      eLpNorm (M (Icc (1 : Real) 2) f) p (powerWeight 2 alpha)
        <= ENNReal.ofReal C * eLpNorm f p (powerWeight 2 alpha) := by
  -- Duoandikoetxea--Vega local proof
  sorry

end Spherical.PowerWeights
\end{lstlisting}

The ASCII spellings in the listing are schematic.  The actual source should use
the project's Unicode notation and exact imported namespace names.

\subsection{The global direct corollary}

\begin{lstlisting}
/-- Full planar circular maximal theorem in the exact missing range. -/
theorem eLpNorm_circularMaximal_powerWeight_le_of_neg
    {p : ENNReal} {alpha : Real}
    (hp_top : p != top) (hp : (2 : ENNReal) < p)
    (halpha_lo : -1 < alpha) (halpha_hi : alpha < 0) :
    exists C : Real, 0 < C and forall f : (Real^2) -> Complex,
      MemLp f p (powerWeight 2 alpha) ->
      MemLp (M (Ioi (0 : Real)) f) p (powerWeight 2 alpha) and
      eLpNorm (M (Ioi (0 : Real)) f) p (powerWeight 2 alpha)
        <= ENNReal.ofReal C * eLpNorm f p (powerWeight 2 alpha) := by
  -- Apply the local theorem, Bourgain, and the FRS globalization lemma.
  sorry
\end{lstlisting}

\subsection{The full planar source range}

Once the proof is in place, expose the full planar theorem at negligible extra
cost.

\begin{lstlisting}
/-- Planar case of Duoandikoetxea--Vega Theorem 10. -/
theorem eLpNorm_circularMaximal_powerWeight_le
    {p : ENNReal} {alpha : Real}
    (hp_top : p != top) (hp : (2 : ENNReal) < p)
    (halpha_lo : -1 < alpha)
    (halpha_hi : alpha < ENNReal.toReal p - 2) :
    exists C : Real, 0 < C and forall f : (Real^2) -> Complex,
      MemLp f p (powerWeight 2 alpha) ->
      MemLp (M (Ioi (0 : Real)) f) p (powerWeight 2 alpha) and
      eLpNorm (M (Ioi (0 : Real)) f) p (powerWeight 2 alpha)
        <= ENNReal.ofReal C * eLpNorm f p (powerWeight 2 alpha) := by
  sorry
\end{lstlisting}

For the current public theorem, the narrower \lean{alpha < 0} corollary has
simpler arithmetic and avoids carrying \lean{ENNReal.toReal p - 2} through the
FRS case split.

\subsection{Set-monotonicity helpers}

These should be proved once near the definition of \lean{M}.

\begin{lstlisting}
theorem restrictedSphericalMaximal_mono
    {d : Nat} {E F : Set Real} (hEF : E subset F)
    (f : (Real^d) -> Complex) (x : Real^d) :
    M E f x <= M F f x := by
  sorry

lemma memLp_restrictedSphericalMaximal_of_le
    {d : Nat} {E F : Set Real} {p : ENNReal} {mu : Measure (Real^d)}
    (hEF : E subset F) {f : (Real^d) -> Complex}
    (hmem : MemLp (M F f) p mu) :
    MemLp (M E f) p mu := by
  sorry

lemma eLpNorm_restrictedSphericalMaximal_mono
    {d : Nat} {E F : Set Real} {p : ENNReal} {mu : Measure (Real^d)}
    (hEF : E subset F) (f : (Real^d) -> Complex) :
    eLpNorm (M E f) p mu <= eLpNorm (M F f) p mu := by
  sorry
\end{lstlisting}

The last two lemmas may follow from existing Mathlib monotonicity theorems once
measurability is supplied.

\subsection{The arithmetic bridge into the FRS branch}

The integration file should expose a small lemma whose statement mirrors the
actual branch assumptions.

\begin{lstlisting}
lemma neg_one_lt_alpha_of_planar_strict_lower
    {E : Set Real} {p alpha : Real}
    (hstrict : legendreAssouadFunction E (p - 2 - alpha) < p - 1) :
    -1 < alpha := by
  have hrho : p - 2 - alpha <=
      legendreAssouadFunction E (p - 2 - alpha) :=
    legendreAssouadFunction_ge_self E (p - 2 - alpha)
  linarith
\end{lstlisting}

The exact theorem name for \lean{legendreAssouadFunction\_ge\_self} should be
matched to the existing FRS formalization.

\subsection{Final $d=2$ branch, schematic Lean proof}

\begin{lstlisting}
by_cases halpha : 0 <= alpha

case pos =>
  exact existing_nonnegative_weight_branch ...

case neg =>
  have halpha_neg : alpha < 0 := lt_of_not_ge halpha
  by_cases hp2 : pReal <= 2

  case pos =>
    exact existing_frs_p_le_two_branch ...

  case neg =>
    have hp2' : 2 < pReal := lt_of_not_ge hp2
    have halpha_lo : -1 < alpha := by
      exact neg_one_lt_alpha_of_planar_strict_lower hstrictLower
    have hbound :=
      eLpNorm_circularMaximal_powerWeight_le_of_neg
        hp_top hp2_ennreal halpha_lo halpha_neg
    exact transfer_full_bound_to_restricted hEpos hbound
\end{lstlisting}

The existing proof already has conversions between the real parameter used in
the admissible-region coordinates and the \lean{ENNReal} exponent used by
\lean{MemLp}.  Reuse those conversions rather than creating a second parameter
bridge inside the new theorem.

\section{Recommended file structure}

All paths in the first six rows are relative to
\lean{LeanSpherical/Codex/Spherical/PowerWeights/}.

\begin{longtable}{>{\raggedright\arraybackslash\ttfamily}p{0.41\textwidth}p{0.52\textwidth}}
\toprule
\normalfont File & Purpose \\
\midrule
\endhead
DuoVega/Basic.lean &
Sign convention, local maximal operator, rational-radius reduction, finite-grid
operators, and set monotonicity. \\
\addlinespace
DuoVega/LittlewoodPaley.lean &
Frequency projections, reconstruction, low-frequency estimate, kernel bounds,
and weighted projection bounds. \\
\addlinespace
DuoVega/Unweighted.lean &
Adapter from the formalized Bourgain theorem to the frequency-localized decay
estimate. \\
\addlinespace
DuoVega/CriticalWeight.lean &
Dyadic spatial annuli, source equation-(3) estimate, off-diagonal kernel decay,
and the $2^{j\varepsilon}$ critical-weight bound. \\
\addlinespace
DuoVega/Interpolation.lean &
Finite-grid vector-valued interpolation, fixed-measure $2$-to-$p$ step, and
Stein--Weiss change-of-measure step. \\
\addlinespace
DuoVega/Theorem.lean &
Local theorem, extension from Schwartz functions, globalization, full source
range, and negative-weight corollary. \\
\addlinespace
PowerWeightTheorem.lean &
One new $d=2$, $p>2$, $\alpha<0$ case; no changes to the other branches. \\
\bottomrule
\end{longtable}

A smaller file split is possible, but the critical-weight proof should not be
mixed with the short FRS integration argument.  Keeping that boundary clear
makes it possible to test the new analytic theorem independently.

\section{Formalization risks and their resolutions}

\begin{longtable}{p{0.28\textwidth}p{0.31\textwidth}p{0.34\textwidth}}
\toprule
Issue & Failure mode & Recommended resolution \\
\midrule
\endhead
Uncountable supremum & Measurability and interpolation do not apply directly. &
Reduce to rational radii for Schwartz functions; interpolate finite grids;
pass monotonically to the countable supremum. \\
\addlinespace
Sublinearity & Riesz--Thorin expects a linear map. &
Use the finite-dimensional vector-valued map
$f\mapsto(\A_tP_jf)_{t\in F}$. \\
\addlinespace
Only final Bourgain theorem exposed & Global boundedness does not imply a
frequency-decaying norm. & Expose the frequency-localized local estimate from
the Bourgain proof as a separate theorem. \\
\addlinespace
Weight singularity at zero & Pointwise density identities fail at $x=0$. &
State density comparisons almost everywhere and remove the singleton using
its nullity. \\
\addlinespace
Negative real powers in \lean{ENNReal} & Simplification rules may create
\lean{top} at zero. & Isolate all power-weight a.e. lemmas in one file. \\
\addlinespace
\lean{p : ENNReal} may equal \lean{top} & The hypothesis \lean{2 < p} alone
does not imply finiteness. & Include \lean{p != top} explicitly or state the
analytic theorem with a real exponent and convert. \\
\addlinespace
Normalization of surface measure & Source constant does not syntactically
match project average. & Prove a scalar relation once and absorb the scalar in
$C$. \\
\addlinespace
Sign $x-t\omega$ versus $x+t\omega$ & Source operator does not unfold to the
project operator. & Prove reflection invariance of normalized sphere measure. \\
\addlinespace
Positive constant in \lean{typeSet} & An operator norm may be returned merely
as a nonnegative real. & Replace $C$ by $\max\{C,1\}$ at the API boundary. \\
\addlinespace
Extension from Schwartz functions & The pointwise supremum is not a continuous
linear operator. & Use sublinearity, density, a Cauchy subsequence, and Fatou,
or reuse an existing maximal-extension lemma. \\
\bottomrule
\end{longtable}

\section{Dependency checklist}

The following table separates assumptions stated in the question from genuinely
new work.

\begin{longtable}{p{0.49\textwidth}p{0.18\textwidth}p{0.25\textwidth}}
\toprule
Component & Status under the question & Action \\
\midrule
\endhead
Bourgain full circular maximal theorem for $p>2$ & Assumed formalized & Reuse. \\
Bourgain frequency-localized decay & Must be checked in the existing API &
Expose if hidden in the proof. \\
FRS Legendre--Assouad inequalities & Assumed formalized & Reuse. \\
FRS strict-region equivalence & Assumed formalized & Reuse. \\
FRS $p\leq2$, $\alpha<0$ upper bound & Assumed formalized & Reuse. \\
FRS nonnegative-weight branch & Assumed formalized & Reuse. \\
FRS local-to-global Lemma 2.1 & Assumed formalized & Reuse. \\
Power-weight and maximal-function set monotonicity & Elementary & Add helper lemmas. \\
Local DV theorem for $p>2$, $-1<\alpha<0$ & Missing & Main new proof. \\
Global planar DV corollary & Missing & Short corollary. \\
$d=2$ case split in the public theorem & Missing & Short integration patch. \\
Necessity, closure, endpoint approximation & Already formalized & No change. \\
\bottomrule
\end{longtable}

\section{Acceptance criteria}

The planar extension should be considered complete only when all of the
following hold.

\begin{enumerate}[leftmargin=8mm]
\item The local negative-weight theorem has no \lean{sorry} and is stated for
      the project's actual \lean{M} and \lean{powerWeight}.
\item The theorem handles every finite exponent $p>2$, not merely a compact
      subrange or integer exponents.
\item The constants are uniform over normalized slices $E_R\subset[1,2]$.
\item The output \lean{MemLp} assertion is proved, not inferred informally from
      a finite norm.
\item The global full-radii theorem is derived through a checked globalization
      lemma.
\item The restricted-radii transfer is pointwise and does not use an equality
      of supremum sets that fails at nonpositive radii.
\item The $d=2$ closure theorem compiles with hypothesis \lean{2 <= d}.
\item No endpoint theorem at $\alpha=-1$ is silently assumed.
\item The final source tree contains no axioms or declarations that merely
      restate the target estimate without proof.
\end{enumerate}

\section{Answer to the central question}

The prior agent's diagnosis is mathematically correct, with one refinement.
The exact direct missing claim is

\[
p>2,
\qquad
-1<\alpha<0,
\qquad
\M:L^p(\R^2,\abs{x}^{\alpha}\,dx)	o
L^p(\R^2,\abs{x}^{\alpha}\,dx).
\]

The strictly minimal analytic claim is the corresponding unit-scale estimate
for $\M_{[1,2]}$, because the FRS globalization lemma is already assumed
formalized.  Both are contained in the planar specialization of
Duoandikoetxea--Vega Theorem~10, whose full weight range is

\[
-1<\alpha<p-2.
\]

For a strict admissible point in the only uncovered branch, the FRS lower
condition forces $\alpha>-1$, while $p>2$ and $\alpha<0$ force
$\alpha<p-2$.  Therefore the new theorem applies exactly where needed.  All
other arguments in the $d\geq3$ proof carry over without analytic changes.

\appendix

\section{Real-exponent variant of the Lean theorem}

It may be easier to formalize the analytic theorem with $p:\R$ and convert to
\lean{ENNReal.ofReal p} only at the API boundary.

\begin{lstlisting}
theorem eLpNorm_localCircularMaximal_powerWeight_le_ofReal
    {p alpha : Real}
    (hp : 2 < p) (halpha_lo : -1 < alpha) (halpha_hi : alpha < 0) :
    exists C : Real, 0 < C and forall f : (Real^2) -> Complex,
      MemLp f (ENNReal.ofReal p) (powerWeight 2 alpha) ->
      MemLp (M (Icc (1 : Real) 2) f) (ENNReal.ofReal p)
        (powerWeight 2 alpha) and
      eLpNorm (M (Icc (1 : Real) 2) f) (ENNReal.ofReal p)
          (powerWeight 2 alpha)
        <= ENNReal.ofReal C *
          eLpNorm f (ENNReal.ofReal p) (powerWeight 2 alpha) := by
  sorry
\end{lstlisting}

Advantages are ordinary real arithmetic in the interpolation parameters and no
separate \lean{p != top} hypothesis.  The disadvantage is one conversion step
when the theorem is inserted into \lean{typeSet}, whose exponent is already
\lean{ENNReal}.  Either interface is sound; the choice should follow the
existing interpolation API.

\section{Direct proof of the restricted transfer}

Suppose the full theorem gives a constant $C>0$ and let
$E\subset(0,\infty)$.  For every $x$,

\[
\sup_{t\in E}\abs{\A_t f(x)}
\leq
\sup_{t>0}\abs{\A_t f(x)}.
\]

After applying \lean{ENNReal.ofReal}, monotonicity remains valid.  Hence

\[
M(E,f,x)\leq M((0,\infty),f,x).
\]

If the right-hand side is in weighted $L^p$, the lattice property of $L^p$
gives \lean{MemLp} for the left-hand side.  Norm monotonicity gives

\[
\operatorname{eLpNorm}(M(E,f))
\leq
\operatorname{eLpNorm}(M((0,\infty),f)).
\]

Combining with the full estimate proves the restricted estimate with exactly
the same constant.  No property of the fractal set $E$ is used in this transfer.

\section{Bibliography}

{\small
\setlength{\parskip}{0pt}
\begin{thebibliography}{99}

\bibitem{Bourgain1986}
J. Bourgain,
\emph{Averages in the plane over convex curves and maximal operators},
J. Analyse Math. 47 (1986), 69--85.

\bibitem{DV1996}
J. Duoandikoetxea and L. Vega,
\emph{Spherical means and weighted inequalities},
J. London Math. Soc. (2) 53 (1996), no. 2, 343--353.
DOI: 10.1112/jlms/53.2.343.

\bibitem{DS2002}
J. Duoandikoetxea and E. Seijo,
\emph{Weighted inequalities for some spherical maximal operators},
Illinois J. Math. 46 (2002), no. 4, 1299--1312.

\bibitem{FRS2026}
M. Fraccaroli, J. Roos, and A. Seeger,
\emph{Power weight inequalities for spherical maximal functions},
arXiv:2602.17613, 2026; published online in Math. Z.,
DOI: 10.1007/s00209-026-04106-4.

\bibitem{Lee2025}
J. Lee,
\emph{The critical weighted inequalities of the spherical maximal function},
arXiv:2309.00068; J. Geom. Anal. 35 (2025), Paper No. 220.

\bibitem{MSS1992}
G. Mockenhaupt, A. Seeger, and C. Sogge,
\emph{Wave front sets, local smoothing and Bourgain's circular maximal theorem},
Ann. of Math. (2) 136 (1992), no. 1, 207--218.

\bibitem{SteinWeiss1975}
E. M. Stein and G. Weiss,
\emph{Introduction to Fourier Analysis on Euclidean Spaces},
Princeton University Press, 1975.

\end{thebibliography}
}

\end{document}
